% ============================ ARQUITECTURA ============================
\chapter{Arquitectura y decisiones de diseño}\label{cap:arquitectura}

\section{Topología}
El laboratorio se compone de tres redes virtuales de libvirt. La red
default (NAT del host) hace de Internet para OPNsense. La red
soc-lab es la LAN interna donde se encuentran el SIEM y las víctimas. Y la red attack aísla a la máquina atacante.

\captura{topologia-fisica}{Topología del laboratorio (modelada en Packet Tracer)}{topologia}

\begin{table}[H]\centering\small
\begin{tabularx}{\linewidth}{l l l X}
  \toprule
  \textbf{Red libvirt} & \textbf{Puente} & \textbf{Subred} & \textbf{Uso} \\
  \midrule
  default & virbr0       & 192.168.122.0/24 & WAN de OPNsense (NAT + DHCP de libvirt) \\
  soc-lab & virbr-soc    & 10.10.10.0/24    & LAN: SIEM y víctimas \\
  attack  & virbr-attack & 10.10.20.0/24    & Red del atacante (Kali) \\
  \bottomrule
\end{tabularx}
\caption{Redes virtuales del laboratorio}\label{tab:redes}
\end{table}

\begin{table}[H]\centering\small
\begin{tabularx}{\linewidth}{l l c c c X}
  \toprule
  \textbf{VM} & \textbf{IP} & \textbf{RAM} & \textbf{vCPU} & \textbf{Disco} & \textbf{Rol} \\
  \midrule
  opnsense       & 10.10.10.1 / 10.10.20.1 & 2 GB & 2 & 20 GB & Router, firewall, DHCP, Suricata \\
  wazuh          & 10.10.10.10             & 8 GB & 4 & 50 GB & SIEM \\
  windows10      & 10.10.10.130 (DHCP)     & 4 GB & 2 & 50 GB & Víctima Windows con Sysmon y agente \\
  ubuntu-victima & 10.10.10.20             & 2 GB & 1 & 15 GB & Víctima Linux con agente \\
  kali           & 10.10.20.10             & 4GB   & 2 & 80GB   & Atacante \\
  host           & 10.10.10.254            & --   & -- & --   & Administración y acceso al dashboard \\
  \bottomrule
\end{tabularx}
\caption{Inventario de máquinas virtuales}\label{tab:inventario}
\end{table}

\section{Flujo de la telemetría}
\begin{itemize}[leftmargin=*,itemsep=0pt]
  \item Endpoint Windows: registro de seguridad y Sysmon $\rightarrow$ agente Wazuh $\rightarrow$ manager (10.10.10.10).
  \item Endpoint Linux: auth.log, syslog $\rightarrow$ agente Wazuh $\rightarrow$ manager.
  \item Red: Suricata en la LAN de OPNsense $\rightarrow$ eve.json $\rightarrow$ agente Wazuh de OPNsense $\rightarrow$ manager.
\end{itemize}


